\chapter{Source Reference}
\label{chap:source}

A module-by-module tour of \verb|src/|. Earlier chapters explained the
\emph{why}; this is the \emph{where} --- the map you use when editing a specific
file. Public items are named so you can grep for them.

\section{\texttt{src/main.rs}}
The binary entry point. Initializes \verb|tracing_subscriber| with an env filter
defaulting to \verb|info,tower_http=info|, loads the config, connects and
initializes the database, assembles \verb|AppState|, builds the router, binds a
\verb|TcpListener|, and calls \verb|axum::serve| with connect-info so peer IPs
reach handlers. It is intentionally thin and is the one module not covered by
unit tests (its logic is exercised through the integration harness instead).

\section{\texttt{src/lib.rs}}
The library root. It:
\begin{itemize}
  \item declares every module (\verb|auth|, \verb|config|, \verb|db|,
        \verb|db_ops|, \verb|error|, \verb|models|, \verb|rate_limit|,
        \verb|routes|, \verb|validate|, \verb|ws|);
  \item embeds \verb|static/| and \verb|templates/| via \verb|include_dir!|;
  \item defines \verb|AppState| (\S\ref{chap:architecture});
  \item defines \verb|build_router|, the single source of truth for the route
        table;
  \item serves \verb|index| (the embedded \verb|index.html|) and
        \verb|serve_static| (embedded assets with a \verb|mime_guess| content
        type), plus an \verb|internal_error| helper that logs and returns a plain
        500.
\end{itemize}

\section{\texttt{src/config.rs}}
Layered configuration; see Chapter~\ref{chap:config}. Public surface:
\verb|Config| (the resolved struct), \verb|Config::load|, and
\verb|Config::for_test|. Private helpers: \verb|read_yaml_or_embedded|,
\verb|project_root|, \verb|env_or|, \verb|env_or_parse|, and the \verb|YamlConfig|
deserialization structs. Carries a \verb|#[cfg(test)]| module for the cascade.

\section{\texttt{src/db.rs}}
Schemas and lifecycle; see Chapter~\ref{chap:database}. Public: the five schema
functions (\verb|users_schema|, \verb|messages_schema|, \verb|attachments_schema|,
\verb|channels_schema|, \verb|channel_members_schema|), \verb|connect|, and
\verb|init_db|. Private: \verb|migrate_messages_table|.

\section{\texttt{src/db\_ops.rs}}
The query toolkit (\S\ref{chap:database}). Public: \verb|open|, \verb|scan_where|,
\verb|append|, \verb|first_string|, \verb|first_large_string|,
\verb|collect_string_column|, \verb|total_rows|, \verb|utf8_row|, \verb|mixed_row|
(with the \verb|Cell| enum), and \verb|rows_to_stored_messages|. Extensively
unit-tested.

\section{\texttt{src/models.rs}}
The serde DTOs. Requests (\verb|Deserialize|): \verb|RegisterRequest|,
\verb|LoginRequest|, \verb|CreateChannelRequest|, \verb|AddChannelMemberRequest|,
\verb|ChannelEncryptedEnvelope|, \verb|ChannelAttachment|,
\verb|ChannelMessagePayload|. Query strings: \verb|TokenQuery|, \verb|SearchQuery|,
\verb|AfterQuery|. Responses (\verb|Serialize|): \verb|TokenResponse|,
\verb|UserListItem|, \verb|UserPublicKeyResponse|, \verb|ChannelInfo|,
\verb|ChannelBrowseItem|, \verb|ChannelMemberKey|, \verb|StoredMessage|,
\verb|AttachmentResponse|. Note the \verb|#[serde(default)]| attributes: channel
\verb|members| default to empty, \verb|after| to \verb|0.0|, and
\verb|message_type| to \verb|"text"| via \verb|default_message_type|.

\section{\texttt{src/auth.rs}}
bcrypt and JWT; see Chapter~\ref{chap:security}. Public: \verb|hash_password|,
\verb|verify_password|, \verb|create_token|, \verb|require_auth|,
\verb|decode_token|. Private: the \verb|Claims| struct and \verb|algorithm_from|.
Unit tests cover the algorithm mapping, password round-trips, token round-trips,
wrong-secret and garbage rejection, and expiry.

\section{\texttt{src/validate.rs}}
The input guards \verb|validate_username|, \verb|validate_id|, and
\verb|validate_after| (\S\ref{chap:security}). \verb|validate_id| and
\verb|validate_after| are marked \verb|#[allow(dead_code)]| at the definition
site but are used by the channel/message handlers; the attribute simply suppresses
a false-positive lint. Unit-tested for accepted and rejected inputs.

\section{\texttt{src/rate\_limit.rs}}
The \verb|RateLimiter| sliding window (\S\ref{chap:security}): \verb|new|,
\verb|check|, and a test-only \verb|clear|. The window (60\,s) and cap (20) are
module constants. Unit-tested for the limit boundary, per-IP isolation, and
window reset.

\section{\texttt{src/error.rs}}
\verb|AppError| and its HTTP mapping (\S\ref{chap:security}). Variants:
\verb|BadRequest(String)|, \verb|Unauthorized(&'static str)|,
\verb|Forbidden(&'static str)|, \verb|NotFound(String)|, \verb|TooManyRequests|,
and \verb|Internal(#[from] anyhow::Error)|. \verb|AppResult<T>| is the handler
return alias. \verb|IntoResponse| renders \verb|{"detail": ...}| with the right
status and masks internals. Async unit tests assert each variant's status/body
and the masking of internal errors.

\section{\texttt{src/ws.rs}}
The \verb|WsConnections| registry (\S\ref{chap:websocket}): \verb|register|,
\verb|unregister|, \verb|send_to|, \verb|send_to_except|, over a
\verb|Mutex<HashMap<String, Vec<(Uuid, ConnTx)>>>|. Unit-tested for fan-out to
multiple connections, the exclusion path, no-op sends, and last-connection
removal.

\section{\texttt{src/routes/}}
Handlers, one module per resource. \verb|mod.rs| declares them.

\subsection*{\texttt{auth.rs}}
\verb|register| and \verb|login| --- rate-limited, validated, bcrypt/JWT-backed
(Chapter~\ref{chap:http-api}).

\subsection*{\texttt{users.rs}}
\verb|list_users| (DM peers derived from the messages table),
\verb|search_users| (rate-limited prefix search, capped at 20), and
\verb|get_public_key|.

\subsection*{\texttt{channels.rs}}
The largest handler module. Public handlers: \verb|create_channel|,
\verb|list_channels|, \verb|browse_channels|, \verb|join_channel|,
\verb|get_channel|, \verb|add_channel_member|, \verb|get_channel_member_keys|.
Private helpers worth knowing: \verb|list_channel_member_usernames|,
\verb|is_channel_member|, \verb|fetch_channel| (returns a private
\verb|ChannelRow|), \verb|channel_info_for|, \verb|build_member_batch| (inserts
all members in one \verb|RecordBatch|), and \verb|now_secs|.

\subsection*{\texttt{messages.rs}}
\verb|get_dm_messages|, \verb|get_channel_messages|, and
\verb|post_channel_message|. The shared engine lives here too:
\verb|store_channel_message| (persist one row per envelope + optional attachment,
returning \verb|ChannelMessageStoreResult|) and \verb|relay_channel_message|
(push one live frame per envelope). Both are called by the REST handler and by the
WebSocket channel handler --- the single point that guarantees the two transports
behave identically. \verb|assert_channel_member| enforces membership.

\subsection*{\texttt{attachments.rs}}
\verb|get_attachment| --- fetch an encrypted blob with a channel-membership check
when the attachment is owned by a channel.

\subsection*{\texttt{ws.rs}}
The WebSocket handler (Chapter~\ref{chap:websocket}): \verb|ws_chat| (upgrade +
token gate), \verb|handle_socket| (registry + writer task + dispatch loop),
\verb|handle_dm| (the double-write and three-way relay), and \verb|handle_channel|
(membership check, then the shared store/relay). \verb|text_msg| wraps a
\verb|serde_json::Value| into a WS text frame.

\section{Cross-Cutting Idioms}
A handful of patterns recur and are worth internalizing before making changes:
\begin{itemize}
  \item \textbf{Thin handlers.} Validate $\to$ authenticate $\to$ authorize $\to$
        call \verb|db_ops|/\verb|ws| $\to$ map errors. Business logic that two
        handlers share is factored out (as with the channel store/relay pair).
  \item \textbf{Timestamps are \texttt{f64} Unix seconds}, produced by a local
        \verb|now_secs()| in each module that needs it.
  \item \textbf{Errors flow through \texttt{AppError}}; use \verb|?| freely and
        let unexpected failures become masked 500s.
  \item \textbf{Never interpolate an unvalidated identifier} into a predicate.
\end{itemize}
